% theory_market_impact.tex — BTC Market Impact and Reflexivity
% Price impact of MSTR accumulation, reflexivity amplification, systemic risk

\section{BTC Market Impact and Reflexivity}\label{sec:market_impact}

This section analyzes the feedback between Strategy's BTC accumulation and
the BTC spot market.  We model the price impact of large-scale purchases,
derive the reflexivity amplification factor that links the equity premium
to BTC price dynamics, and assess systemic risk implications as the
treasury vehicle model scales.

%% ====================================================================
\subsection{Price Impact of BTC Accumulation}\label{subsec:price_impact}

Strategy is among the largest single-entity BTC buyers, with purchases
that can represent a material fraction of daily trading volume.  Following
\citet{kyle1985continuous}, we incorporate a permanent price impact into
the BTC dynamics.

\begin{definition}[Market Impact Model]\label{def:impact_model}
The BTC price process with endogenous market impact is
\begin{equation}\label{eq:btc_impact}
  \frac{dS_t}{S_t}
  = \mu_S\, dt + \sigma_S\, dW_t^S
    + \eta \cdot \frac{dH_t}{V_t},
\end{equation}
where $V_t$ is the daily BTC spot trading volume (in BTC), $dH_t$ is
Strategy's BTC purchase flow, and $\eta > 0$ is the Kyle lambda
(permanent price impact coefficient), measuring the fractional price
change per unit of order flow normalized by volume.
\end{definition}

\begin{remark}
The impact term $\eta \cdot dH_t / V_t$ captures two channels:
(i) \emph{information impact}---the market infers bullish private
information from Strategy's purchases; and (ii) \emph{mechanical
impact}---large buy orders move the order book.  For a treasury
vehicle with publicly announced purchases, channel (i) dominates
through the signaling effect of announced acquisitions.
\end{remark}

%% ====================================================================
\subsection{Estimation of Price Impact}\label{subsec:eta_estimate}

\begin{proposition}[Price Impact Estimation]\label{prop:eta_est}
The price impact coefficient $\eta$ can be estimated from an event study
around Strategy's BTC purchase announcements.  Using the calibrated
parameters from Section~\ref{sec:methodology}:
\begin{itemize}
  \item Average purchase size: $\overline{\Delta H} \approx 6{,}887$ BTC
    per event (from the empirical jump size distribution $\mathcal{D}_H$).
  \item Average daily BTC spot volume:
    $\bar{V} \approx 250{,}000\text{--}400{,}000$ BTC
    (approximately \$20--30B at $S \approx \$80{,}000$).
  \item Order flow fraction per event:
    $\overline{\Delta H} / \bar{V} \approx 0.017\text{--}0.028$.
\end{itemize}
\end{proposition}

\begin{definition}[Event Study Estimator]\label{def:event_study}
Let $\{t_k\}_{k=1}^K$ denote the $K$ purchase announcement dates with
purchase sizes $\{\Delta H_k\}$.  The abnormal return around event $k$ is
\begin{equation}\label{eq:abnormal_return}
  \mathrm{AR}_k
  = \frac{S_{t_k + 1} - S_{t_k}}{S_{t_k}}
    - \hat{\mu}_S \cdot \Delta t,
\end{equation}
where $\hat{\mu}_S$ is the estimated daily BTC drift.  The price impact
coefficient is estimated via OLS:
\begin{equation}\label{eq:eta_ols}
  \mathrm{AR}_k = \eta \cdot \frac{\Delta H_k}{V_{t_k}} + \varepsilon_k,
  \qquad \varepsilon_k \sim (0, \sigma_\varepsilon^2).
\end{equation}
\end{definition}

\begin{proposition}[Bounds on $\eta$]\label{prop:eta_bounds}
Under the \citet{kyle1985continuous} framework with normally distributed
order flow, the theoretical price impact satisfies
\begin{equation}\label{eq:kyle_lambda}
  \eta = \frac{\sigma_S}{\sqrt{V_t}} \cdot \frac{1}{\sigma_u},
\end{equation}
where $\sigma_u$ is the standard deviation of uninformed order flow.
For BTC markets with $\sigma_S \approx 0.60$ (annualized) and
$V_t \approx 300{,}000$ BTC/day, this yields
\begin{equation}\label{eq:eta_calibrated}
  \eta \approx 0.5\text{--}2.0,
\end{equation}
implying that a purchase of $\overline{\Delta H} / V_t \approx 2\%$ of
daily volume generates a $1\text{--}4\%$ price impact.
\end{proposition}

\begin{remark}
Empirical estimates from cryptocurrency market microstructure literature
suggest $\eta$ in the range of $0.5$--$3.0$ for large-cap tokens,
consistent with our calibration range.  The higher end reflects the
signaling premium attached to Strategy's purchases: the market interprets
MSTR buying as a bullish signal, amplifying the mechanical impact.
\end{remark}

%% ====================================================================
\subsection{Reflexivity Amplification}\label{subsec:reflexivity}

The market impact channel creates a reflexive feedback loop:
Strategy buys BTC $\to$ BTC price rises $\to$ NAV rises $\to$
premium rises $\to$ more issuance capacity $\to$ more BTC buying.
We formalize this loop and derive the amplification factor.

\begin{definition}[Feedback Chain]\label{def:feedback_chain}
The reflexivity loop consists of four marginal sensitivities:
\begin{enumerate}
  \item \textbf{Price impact:} $\partial S / \partial H = \eta S / V$
    --- BTC price response to purchases.
  \item \textbf{NAV sensitivity:} $\partial \mathrm{NAV} / \partial S
    = H$ --- NAV response to BTC price.
  \item \textbf{Premium response:} $\partial \pi / \partial \mathrm{NAV}$
    --- premium response to NAV changes, determined by the OU dynamics
    and the equity identity.  From $E = e^\pi \cdot \mathrm{NAV}$ with
    $E$ responding to NAV, this is approximately
    $\partial \pi / \partial \mathrm{NAV}
    = (1 - e^{-\pi}) / \mathrm{NAV}$ when the market partially
    passes through NAV changes to premium.
  \item \textbf{Issuance response:}
    $\partial H^{\mathrm{new}} / \partial \pi$ --- additional BTC
    acquired in response to premium changes, from the holding dynamics
    $dH_t = \alpha \pi_t^+ dt + dH_t^{\mathrm{jump}}$.  This is
    approximately $\alpha + \lambda_M \partial(\mathbb{E}[\Delta H])
    / \partial \pi$.
\end{enumerate}
\end{definition}

\begin{theorem}[Reflexivity Amplification Factor]\label{thm:reflexivity}
Define the single-loop gain as the product of marginal sensitivities
around the feedback chain:
\begin{equation}\label{eq:loop_gain}
  \mathcal{G}_t
  = \eta \cdot \frac{S_t}{V_t}
    \cdot H_t
    \cdot \frac{\partial \pi}{\partial \mathrm{NAV}}\bigg|_t
    \cdot \frac{\partial H^{\mathrm{new}}}{\partial \pi}\bigg|_t.
\end{equation}
The reflexivity amplification factor is
\begin{equation}\label{eq:amplification}
  \mathcal{R}_t = \frac{1}{1 - \mathcal{G}_t},
\end{equation}
provided $\mathcal{G}_t < 1$.  The effective BTC price response to an
exogenous demand shock $\Delta H^{\mathrm{exog}}$ is
\begin{equation}\label{eq:amplified_impact}
  \frac{\Delta S}{S}
  = \mathcal{R}_t \cdot \eta \cdot \frac{\Delta H^{\mathrm{exog}}}{V_t}.
\end{equation}
\end{theorem}

\begin{proof}
Consider a purchase of $\Delta H_0$ BTC.  The first-round price impact is
$\Delta S_0 / S = \eta \cdot \Delta H_0 / V$.  This raises NAV by
$H \cdot \Delta S_0$, which raises the premium by
$(\partial\pi/\partial\mathrm{NAV}) \cdot H \cdot \Delta S_0$, which
induces additional issuance of
$(\partial H^{\mathrm{new}}/\partial\pi) \cdot \Delta\pi$ BTC.  The
additional purchase creates a second-round impact, and so on.  The total
impact is the geometric series:
\[
  \frac{\Delta S^{\mathrm{total}}}{S}
  = \eta \frac{\Delta H_0}{V}
    \sum_{n=0}^{\infty} \mathcal{G}^n
  = \frac{\eta \cdot \Delta H_0 / V}{1 - \mathcal{G}},
\]
converging when $\mathcal{G} < 1$.
\end{proof}

\begin{proposition}[Instability Condition]\label{prop:instability}
The reflexivity loop becomes unstable ($\mathcal{R}_t \to \infty$) when
$\mathcal{G}_t \to 1$.  This occurs when:
\begin{equation}\label{eq:instability}
  \eta \cdot \frac{H_t S_t}{V_t}
  \cdot \frac{\partial \pi}{\partial \mathrm{NAV}}\bigg|_t
  \cdot \frac{\partial H^{\mathrm{new}}}{\partial \pi}\bigg|_t
  \;\ge\; 1.
\end{equation}
The loop gain $\mathcal{G}_t$ is increasing in:
\begin{enumerate}
  \item $\eta$: higher price impact (thinner markets);
  \item $H_t S_t / V_t$: Strategy's market footprint (holdings relative
    to trading volume);
  \item $\partial\pi / \partial\mathrm{NAV}$: premium sensitivity
    (higher when the market is momentum-driven);
  \item $\partial H^{\mathrm{new}} / \partial\pi$: issuance aggressiveness
    (higher $\alpha$ or $\lambda_M$).
\end{enumerate}
\end{proposition}

\begin{corollary}[Calibrated Loop Gain]\label{cor:loop_gain}
Using the calibrated parameters:
\begin{align}
  \eta &\approx 1.0, &
  H_t &\approx 500{,}000\;\text{BTC}, \notag \\
  S_t &\approx \$80{,}000, &
  V_t &\approx 300{,}000\;\text{BTC/day}, \notag \\
  \frac{\partial\pi}{\partial\mathrm{NAV}} &\approx 0.02, &
  \frac{\partial H^{\mathrm{new}}}{\partial\pi} &\approx 5{,}000
    \;\text{BTC per unit}\;\pi,
\end{align}
the loop gain is approximately
\begin{equation}\label{eq:G_calibrated}
  \mathcal{G} \approx 1.0 \times \frac{500{,}000 \times 80{,}000}
  {300{,}000 \times 80{,}000} \times 0.02 \times 5{,}000
  \approx 1.0 \times 1.67 \times 100
  = 167,
\end{equation}
which far exceeds unity.  However, this linear approximation breaks down
because the marginal sensitivities are themselves functions of the state
and exhibit strong diminishing returns.  In particular:
\begin{itemize}
  \item The premium response $\partial\pi/\partial\mathrm{NAV}$ declines
    rapidly as $\pi$ rises (mean reversion dampens further increases).
  \item The issuance response $\partial H^{\mathrm{new}}/\partial\pi$
    is bounded by ATM capacity and regulatory constraints.
  \item The price impact $\eta$ may decrease at larger order sizes
    as the market adapts.
\end{itemize}
The effective (nonlinear) loop gain is
\begin{equation}\label{eq:G_effective}
  \mathcal{G}_t^{\mathrm{eff}}
  = \eta^{\mathrm{eff}}(\Delta H)
    \cdot \frac{H_t S_t}{V_t}
    \cdot \left(\frac{\partial\pi}{\partial\mathrm{NAV}}\right)^{\mathrm{eff}}
    \cdot \left(\frac{\partial H^{\mathrm{new}}}
    {\partial\pi}\right)^{\mathrm{eff}},
\end{equation}
where the ``eff'' superscripts denote the effective marginal sensitivities
after accounting for saturation.  Numerical simulation suggests
$\mathcal{G}^{\mathrm{eff}} \in [0.3, 0.8]$ under typical conditions,
yielding $\mathcal{R} \in [1.4, 5.0]$: meaningful amplification but
bounded away from instability.
\end{corollary}

\begin{remark}
The fact that $\mathcal{G}^{\mathrm{eff}} < 1$ while the linearized
$\mathcal{G} \gg 1$ reflects the strong nonlinear damping in the system.
The OU mean reversion of the premium ($\kappa > 0$) is the primary
stabilizing force: without it, the system would indeed be unstable,
as in \citet{soros1987alchemy}'s notion of unchecked reflexivity.
\end{remark}

%% ====================================================================
\subsection{Impact on BTC Volatility}\label{subsec:vol_impact}

\begin{proposition}[Volatility Amplification]\label{prop:vol_amp}
Under the market impact model~\eqref{eq:btc_impact}, the effective BTC
volatility incorporating Strategy's feedback is
\begin{equation}\label{eq:vol_amplified}
  \sigma_S^{\mathrm{eff}}
  = \sigma_S \cdot \sqrt{\mathcal{R}_t^2
    + 2\mathcal{R}_t(\mathcal{R}_t - 1)
    \cdot \rho_{\mathrm{flow}} \cdot \frac{\sigma_H}{\sigma_S}
    \cdot \frac{S_t}{V_t}},
\end{equation}
where $\sigma_H$ is the volatility of Strategy's purchase flow and
$\rho_{\mathrm{flow}}$ is the correlation between BTC returns and
Strategy's purchase decisions.  When $\mathcal{R}_t > 1$:
\begin{equation}\label{eq:vol_increase}
  \sigma_S^{\mathrm{eff}} > \sigma_S.
\end{equation}
Strategy's presence in the market unambiguously increases BTC volatility
when the reflexivity amplification factor exceeds unity.
\end{proposition}

\begin{proof}
From~\eqref{eq:btc_impact}, the instantaneous variance of $dS/S$ is
\begin{align*}
  \mathrm{Var}\!\left(\frac{dS}{S}\right)
  &= \sigma_S^2\,dt
    + \eta^2 \frac{\mathrm{Var}(dH)}{V^2}
    + 2\eta \frac{\sigma_S \mathrm{Cov}(dW^S, dH)}{V}.
\end{align*}
Incorporating the reflexive response $dH = dH^{\mathrm{exog}}
+ (\partial H/\partial S)\, dS$ and solving for the effective variance
yields the stated expression after substituting
$\mathcal{R} = 1/(1-\mathcal{G})$.
\end{proof}

%% ====================================================================
\subsection{Systemic Risk: Fire Sales and Contagion}\label{subsec:systemic}

As Strategy accumulates a growing share of the total BTC supply, and as
other firms adopt the treasury vehicle model, the systemic implications
become material.  We analyze forced liquidation risk in the framework of
\citet{shleifer1997limits} and \citet{shleifer2011fire}.

\begin{assumption}[Treasury Vehicle Ecosystem]\label{ass:ecosystem}
Suppose $m$ firms adopt the BTC treasury model, with aggregate holdings
$\sum_{i=1}^m H_t^{(i)} = \bar{H}_t$.  Let $\bar{H}_t / \bar{S}$
represent the fraction of circulating BTC supply held by treasury
vehicles, where $\bar{S}$ is the total circulating supply
($\approx 19.8$ million BTC).
\end{assumption}

\begin{definition}[Fire Sale Discount]\label{def:fire_sale}
If firm $i$ is forced to liquidate $\Delta H^{(i)}$ BTC (due to
debt maturity, margin calls, or dividend obligations), the fire sale
discount is
\begin{equation}\label{eq:fire_discount}
  \delta_{\mathrm{fire}}
  = \eta \cdot \frac{\Delta H^{(i)}}{V_t}
    + \eta_{\mathrm{crowd}} \cdot \frac{\sum_{j \ne i}
    \Delta H_{\mathrm{forced}}^{(j)}}{V_t},
\end{equation}
where $\eta_{\mathrm{crowd}}$ captures the contagion effect: firm $i$'s
liquidation depresses BTC price, pushing other treasury vehicles toward
their distress boundaries, triggering further forced sales.
\end{definition}

\begin{theorem}[Liquidation Cascade]\label{thm:cascade}
Under Assumption~\ref{ass:ecosystem}, a liquidation cascade occurs when
an initial price shock $\Delta S / S < 0$ triggers a chain of forced sales.
Let $\ell^{(i)}_t = L_t^{(i)} / A_t^{(i)}$ be firm $i$'s leverage and
$\bar{\ell}^{(i)}$ its distress threshold.  The cascade condition is:
\begin{equation}\label{eq:cascade_condition}
  \exists\; S' < S_t \text{ s.t. }
  \sum_{i:\, \ell^{(i)}(S') \ge \bar{\ell}^{(i)}}
  H^{(i)} \cdot \eta / V_{S'}
  \;>\; \frac{S_t - S'}{S_t},
\end{equation}
i.e., the aggregate price impact of forced liquidations exceeds the
initial price decline that triggered them, creating a self-reinforcing
downward spiral.
\end{theorem}

\begin{proof}
Firm $i$ enters distress when $A^{(i)} = H^{(i)} S \le
(1+\varepsilon) L^{(i)}$, i.e.\ when
$S \le (1+\varepsilon) L^{(i)} / H^{(i)} \equiv S^{*(i)}$.
Starting from $S_t$, a decline to $S'$ triggers liquidation by all
firms with $S^{*(i)} \ge S'$.  The aggregate selling pressure
$\sum_{i \in \mathcal{D}(S')} H^{(i)}$ creates price impact
$\eta \sum_{i \in \mathcal{D}} H^{(i)} / V_{S'}$.  If this impact
exceeds $|S' - S_t|/S_t$, the price must fall further, triggering
additional firms, and the cascade propagates.
\end{proof}

\begin{proposition}[Current Systemic Footprint]\label{prop:footprint}
As of early 2026, Strategy holds approximately $500{,}000$ BTC
($\approx 2.5\%$ of circulating supply).  The systemic risk metrics are:
\begin{enumerate}
  \item \textbf{Market share of daily volume:}
    Average purchase $\overline{\Delta H} / \bar{V}
    \approx 6{,}887 / 300{,}000 \approx 2.3\%$ per event.
  \item \textbf{Forced liquidation impact:}
    Full liquidation of $500{,}000$ BTC at $\eta = 1$ over 30 days
    ($V = 300{,}000 \times 30 = 9{,}000{,}000$ BTC-equivalent volume)
    implies $\delta_{\mathrm{fire}} \approx 500{,}000 / 9{,}000{,}000
    \approx 5.6\%$ price impact, before cascade effects.
  \item \textbf{With $m = 10$ imitator firms} holding an aggregate
    additional $200{,}000$ BTC, the forced-liquidation scenario
    becomes significantly more severe:
    $\delta_{\mathrm{fire}}^{\mathrm{agg}} \approx 7.8\%$ direct impact,
    with cascade multiplier $\mathcal{R}_{\mathrm{crisis}}$ potentially
    exceeding $3\times$ due to correlated leverage and concentrated
    ownership.
\end{enumerate}
\end{proposition}

\begin{remark}[Shleifer--Vishny Fire Sale Channel]
The BTC treasury model creates a classic
\citet{shleifer1997limits} ``limits of arbitrage'' scenario: during
distress, the natural buyers (other treasury vehicles) face the same
adverse conditions and cannot absorb the supply.  Unlike equities, where
diverse investor bases provide liquidity, BTC treasury vehicles share
correlated funding and valuation mechanisms, concentrating liquidation
risk.  The absence of a lender of last resort for BTC exacerbates this
fragility.
\end{remark}

\begin{proposition}[Volatility Regime Shift]\label{prop:vol_regime}
As the fraction of BTC held by treasury vehicles
$\omega_t = \bar{H}_t / \bar{S}$ increases, BTC volatility transitions
between regimes:
\begin{equation}\label{eq:vol_regimes}
  \sigma_S^{\mathrm{eff}}(\omega) \approx
  \begin{cases}
    \sigma_S & \text{if } \omega < \omega_{\mathrm{low}} \approx 0.01, \\
    \sigma_S \cdot (1 + \beta_\omega \omega) & \text{if }
      \omega_{\mathrm{low}} \le \omega \le \omega_{\mathrm{high}}, \\
    \sigma_S \cdot \mathcal{R}(\omega) & \text{if }
      \omega > \omega_{\mathrm{high}} \approx 0.10,
  \end{cases}
\end{equation}
where $\beta_\omega > 0$ captures the linear amplification regime and
$\mathcal{R}(\omega)$ is the full reflexivity amplification factor, which
grows nonlinearly in $\omega$.
\end{proposition}

\begin{remark}
At the current $\omega \approx 0.025$ (Strategy alone), the system is
in the linear amplification regime.  If the treasury model is widely
adopted and $\omega$ exceeds $0.10$, the nonlinear regime would introduce
meaningful excess volatility and tail risk to BTC markets.  This creates
a paradox: the more successful the treasury vehicle model, the more
volatile the underlying asset becomes, which in turn increases the
distress probability for all participants---a collective action problem
in the spirit of \citet{brunnermeier2009market}.
\end{remark}
